\documentclass[a4paper,twocolumn,landscape]{article}
\usepackage{geometry}
\usepackage{concmath}
\usepackage[OT1]{fontenc}
\usepackage{graphicx}
\usepackage{amsfonts}
\usepackage{amsmath}
\usepackage{calrsfs}
\usepackage{color}
\usepackage[dvipsnames]{xcolor}
\usepackage[many]{tcolorbox}
\usepackage{setspace}              
\usepackage{multicol}  
\geometry{margin=0.5in}
\definecolor{green}{RGB}{52, 235, 168}
\definecolor{violet}{RGB}{165, 83, 237}
\definecolor{lemon}{RGB}{170, 237, 83}
\definecolor{red}{RGB}{240, 125, 84}
\definecolor{blue}{RGB}{105, 255, 242}
\begin{document}

\noindent
\colorbox{green}{Exercice 2 :} \\

\noindent
a/ \begin{math}
    \frac{1}{1+j(\frac{\omega}{R}L) - \frac{1}{R\omega C}} = \frac{H_0}{1 + jQ(\frac{\omega}{\omega_0} - \frac{\omega_0}{\omega})}\\
    H_0 = 1\\
    \begin{cases}
        \frac{L}{R} = \frac{Q}{\omega_0}.1\\
        \frac{1}{RC} = \omega_0  Q . 2\\
    \end{cases}
    \begin{cases}
        \frac{2}{1} \to \frac{1}{LC} = \omega_0^2\\
        \Rightarrow \omega_0 = \frac{1}{\sqrt{LC}}\\
        Q = \frac{1}{R}\sqrt{\frac{L}{C}}
    \end{cases}\\
    \omega_0 = 5000rad.s^{-1}\\
    Q = 100 \\
    \omega_0 = 3\omega
\end{math}

\noindent b/ 
\begin{equation*}
    U_s(t) = -\frac{8A}{\pi^2} \sum_{k=0}^{+\infty} \frac{|H(j(2k+1)\omega)|}{(2k+1)^2}cos((2k+1)\omega t + arg(\underline{H}(j(2k+1))\omega))
\end{equation*}

\noindent
c/ \begin{align*}
    A_k &= -\frac{8A}{\pi^2(2k+1)^2} \times \frac{1}{\sqrt{1+Q^2((2k+1)\frac{\omega}{\omega_0}-\frac{\omega_0}{(2k+1)\omega})}}\\
    \varphi_k &= -arctan(Q(\frac{(2k+1)\omega}{\omega_0} - \frac{\omega_0}{(2k+1)\omega}))
\end{align*}

\noindent
d/ \begin{eqnarray*}
    A_0 \simeq 0,03 V\\
    A_1 \simeq  0.90V\\
    A_2 \simeq  0.0V\\
    \varphi_0 \simeq  -1.58 rad\\
    \varphi_1 \simeq  rad\\
    \varphi_2 \simeq  1.58 rad
\end{eqnarray*}
\noindent
e/ \(u_s \simeq 0.90cos(3\omega t)\)\\

\noindent
\colorbox{green}{Exercice 3}\\
La réponse à un signal continu est un signal continu.\\
Pas de discontinuité lors de la discontinuité de l'entrée\\
Donc passe-bas. \\
\begin{equation*}
    H(\omega) = \frac{H_0}{1+\frac{j}{Q}\frac{\omega}{\omega_0}-\frac{\omega_0^2}{\omega^2}}
\end{equation*} 

\begin{align*}
    H_0 &= 2\\
    Q &\simeq 5 \\
    u_p &= \omega\sqrt{1-\frac{1}{2Q^2}}\\
    &\simeq \omega_0 
\end{align*}

\begin{align*}
    \underline{s} (1 + \frac{j\omega}{Q\omega_0} + \frac{(j\omega)^2}{\omega_0^2}) &= H_0 \underline{e}\\
    s(t) + \frac{1}{Q\omega_0}\frac{ds}{dt} + \frac{1}{\omega_0^2}\frac{d^2s}{dt^2} &= H_0e(t) (= 1, t > 0)\\
    \frac{d^2s}{dt^2} + \frac{\omega_0}{Q}\frac{ds}{dt} + \omega_0^2 s(t) &= \omega_0^2H_0\\
    \Delta = \frac{\omega_0^2}{Q^2} - 4\omega_0^2<0\\
    r_{1,2} = -\frac{\omega_0}{2Q} \pm j\omega_0 \sqrt{1-\frac{1}{4Q^2}}\\
    s(t) = ae^{r_1t} + be^{r_2 t}\\
    s(t) = e^{-\frac{\omega_0}{2Q}t}(ae^{j\omega_pt} + be^{-j\omega_pt}) + H_0 - Acos(\omega_p t + \varphi)
\end{align*}
\begin{equation*}
    s'(t) = H_0 + Ae^{-\frac{\omega_0}{2Q}t} cos(\omega_p t + \varphi)
\end{equation*}
\(t_1\) : premier max de s(t)\\
\(t_2\) : 2e --\\
\begin{align*}
    s(t_1) = 2 + Ae^{-\frac{\omega_0}{2Q}t_1} = 3.5\\
    s(t_2) =2 + Ae^{-\frac{\omega_0}{2Q}(t_1 + T_0)} = 2.8
\end{align*}
\begin{math}
    \begin{cases}
        Ae^{-\frac{\omega_0}{2Q}t_1} = 1.5\\
        Ae^{-\frac{\omega_0}{2Q}(t_1+T_0)} = 0.8
    \end{cases}\\
    \frac{1.5}{0.8} = e^{\frac{\omega_0 T_0}{2Q}} = e^{\frac{\pi}{Q}}\\
    \frac{\pi}{Q} = ln(\frac{1.5}{0.8})\\
    Q = 5
\end{math}

\noindent \colorbox{green}{Exercice 4}\\
a/
\begin{align*}
    e &= u_R + u_C\\
    &= R(i_C + i_R) + u_C\\
    &= RC \frac{du_C}{dt} + u_c + u_c\\
    &= \frac{du}{dt} + \frac{2}{RC}u = \pm \frac{E}{RC}
\end{align*}

c/ \begin{align*}
    j\omega \underline{U} + \frac{2\underline{u}}{RC} = \frac{\underline{e}}{RC}\\
    \frac{\underline{u}}{\underline{e}} = \frac{1}{RC(j\omega + \frac{2}{RC})} \\
    \underline{H}(\omega) = \frac{1}{jRC\omega + 2}\\
    \underline{H} = \frac{\frac{1}{2}}{1 + j\frac{RC\omega}{2}}\\
    = \frac{H_0}{1 + j\frac{\omega}{\omega_c}} \begin{cases}
        H_0 = \frac{1}{2}\\
        \omega_c = \frac{2}{RC} = 2 \omega_0
    \end{cases}
\end{align*}
\begin{itemize}
    \item \(\omega <<< \omega_c\) \begin{math}
        \begin{cases}
            G_{dB} = 20log(\frac{1}{2})\\
            = -6\\
            \varphi = 0
        \end{cases}
    \end{math}
    \item \(\omega >> \omega_c\)\begin{align*}
        G_{dB} = 20log(\frac{H_0\omega_c}{\omega})\\
        &= -6 + 20log(\frac{\omega_c}{\omega})\\
        \varphi = -\frac{\pi}{2}
    \end{align*}
\end{itemize}

\noindent
\colorbox{green}{Exercice 5}\\
1/ \begin{eqnarray*}
    A_0 = \alpha E \\
    E_{eff} = \sqrt{\alpha} E
\end{eqnarray*}
2/ \begin{align*}
    A_1 &= \frac{|sin\alpha\pi|}{\pi} 2 E\\
    &= 0.0100.2E\\
    A_{17} &= \frac{|sin17\alpha\pi|}{17\pi} 2E\\
    &= 0.00953\times 2E\\
    \frac{A_1-A_{17}}{A_1} &= \frac{0.00047}{0.01000} = 0.047
\end{align*}
3/ On observe le gain du filtre en fonction de f\\

\colorbox{green}{Exercice 6}\\
\begin{equation*}
    \underline{H}(j\omega) = \frac{H_0}{1+jQ(\frac{\omega}{\omega_0}-\frac{\omega_0}{\omega})}
\end{equation*}
\(H_0 = 1\\
f_0 = 1000Hz\)\\
Pour \(\omega >> \omega_0\) \begin{eqnarray*}
    \underline{H}(j\omega) \approx \frac{H_0}{jQ\frac{\omega}{\omega_0}}\\
    G = \frac{1}{Q\frac{\omega}{\omega_0}}\\
    G_{dB} = -20 log(Q\frac{\omega}{\omega_0})\\
    -20 = -20logQ\\
    Q = 10\\
    Q = \frac{1}{R}\sqrt{\frac{L}{C}}
\end{eqnarray*}




\end{document}